\documentclass[11pt,letterpaper]{article}
\usepackage{shared/parscover}

% ============================================================================
% Threshold Signatures for Pars Network Infrastructure
% Pars Network Technical Whitepaper PNP-010
% ============================================================================

\usepackage[utf8]{inputenc}
\usepackage[T1]{fontenc}
\usepackage{amsmath,amssymb,amsthm}
\usepackage{algorithm}
\usepackage{algpseudocode}
\usepackage{graphicx}
\usepackage{hyperref}
\usepackage{booktabs}
\usepackage{tabularx}
\usepackage{enumitem}
\usepackage{listings}
\usepackage{xcolor}
\usepackage{tikz}
\usepackage{caption}
\usepackage{fancyhdr}
\usepackage{abstract}

\geometry{margin=1in}

\hypersetup{
    colorlinks=true,
    linkcolor=blue!70!black,
    citecolor=green!50!black,
    urlcolor=blue!60!black
}

\lstset{
    basicstyle=\ttfamily\small,
    breaklines=true,
    frame=single,
    backgroundcolor=\color{gray!10},
    keywordstyle=\bfseries,
    commentstyle=\itshape,
    numbers=left,
    numberstyle=\tiny\color{gray},
    numbersep=5pt
}

\usetikzlibrary{arrows.meta,positioning,shapes.geometric,fit,calc}

\newtheorem{definition}{Definition}[section]
\newtheorem{theorem}{Theorem}[section]
\newtheorem{lemma}{Lemma}[section]
\newtheorem{proposition}{Proposition}[section]
\newtheorem{corollary}{Corollary}[section]

\pagestyle{fancy}
\fancyhf{}
\fancyhead[L]{\small Cyrus Pahlavi, Pars Team}
\fancyhead[R]{\small Threshold Signing}
\fancyfoot[C]{\thepage}

% ============================================================================
\title{%
    \textbf{Threshold Signatures for\\
    Pars Network Infrastructure}\\[0.5em]
    \large Pars Network Technical Whitepaper PNP-010
}

\author{%
    Zach Kelling\\
    Pars DAO\\
    \texttt{research@pars.dev}\\[0.5em]
    \small Version 1.0
}

\date{2026}

% ============================================================================
\begin{document}
\parscoverpage

\thispagestyle{fancy}

% ============================================================================
\begin{abstract}
\noindent
We present a threshold signature framework for the Pars Network that
eliminates single points of failure in key management across three
critical subsystems: guardian-based social recovery wallets, multi-party
governance execution, and MPC bridge signing for diaspora remittances.
We instantiate FROST for Ed25519 threshold signatures on the Pars session
layer, and CGGMP21 for ECDSA threshold signatures on the Ethereum-facing
diaspora bridge.  We formalise the security properties under the
algebraic one-more discrete logarithm assumption, describe key resharing
protocols for guardian rotation, and present a coercion-resistant
signing variant where panic credentials trigger decoy signature
generation.  The framework integrates with Pars social recovery
(guardian networks), shadow governance (anonymous execution), and the
diaspora bridge (cross-chain remittances), providing threshold security
without introducing trusted third parties.
\end{abstract}

\vspace{1em}
\noindent\textbf{Keywords:} threshold signatures, FROST, CGGMP21,
social recovery, MPC bridge, guardian networks, diaspora remittances,
key resharing

\tableofcontents
\newpage

% ============================================================================
\section{Introduction}
\label{sec:intro}

Centralised key custody fails the diaspora threat model.  A single
custodian can be coerced, a single device can be seized, and a single
seed phrase can be extracted under duress.  Threshold cryptography
distributes trust across $n$ parties such that any $t$-of-$n$ can
cooperate to produce a valid signature, while $t-1$ colluding parties
learn nothing about the secret key.

Three Pars subsystems require threshold signing:

\begin{enumerate}
    \item \textbf{Social recovery.}  Guardian networks recover user
    accounts; $t$-of-$n$ guardians must cooperate to authorise key
    rotation without any single guardian holding the complete
    key~\cite{pars_social_recovery}.
    \item \textbf{Governance execution.}  Passed governance proposals
    are executed by a multi-sig committee; threshold signing prevents
    any single committee member from unilaterally executing or blocking
    proposals~\cite{pars_shadow_governance}.
    \item \textbf{Diaspora bridge.}  Cross-chain remittances between
    Pars L2 and Ethereum require bridge signatures; a threshold signing
    committee replaces the single-key bridge operator~\cite{pars_diaspora_bridge}.
\end{enumerate}

\subsection{Choice of Schemes}

\begin{itemize}
    \item \textbf{FROST (Flexible Round-Optimised Schnorr Threshold
    signatures)}~\cite{komlo2020frost} for Ed25519: used on the Pars
    session layer and internal governance, where Schnorr signatures are
    native.
    \item \textbf{CGGMP21}~\cite{canetti2021cggmp} for ECDSA: required
    for Ethereum-compatible bridge signatures on secp256k1.
\end{itemize}

\subsection{Contributions}

\begin{itemize}
    \item Concrete FROST instantiation for Pars guardian networks with
    key resharing for guardian rotation.
    \item CGGMP21 deployment for the diaspora bridge with identifiable
    abort and proactive refresh.
    \item Coercion-resistant threshold signing: panic-triggered decoy
    signatures that satisfy the coercion model of PNP-005.
    \item Performance benchmarks on consumer hardware and mobile devices
    used by diaspora communities.
\end{itemize}

% ============================================================================
\section{Preliminaries}
\label{sec:prelim}

\subsection{Notation}

Let $\mathbb{G}$ be a prime-order group of order $q$ with generator $G$.
For Ed25519, $\mathbb{G}$ is the prime-order subgroup of Curve25519 with
$q = 2^{252} + 27742317777372353535851937790883648493$.  For secp256k1,
$q$ is the curve order.

\begin{definition}[$t$-of-$n$ Secret Sharing]
    A $(t, n)$ Shamir secret sharing of $s \in \mathbb{Z}_q$ distributes
    shares $s_1, \ldots, s_n$ such that any $t$ shares determine $s$ via
    Lagrange interpolation, and any $t-1$ shares reveal no information
    about $s$.
\end{definition}

\begin{definition}[Threshold Signature]
    A $(t, n)$ threshold signature scheme consists of:
    \begin{enumerate}[label=(\roman*)]
        \item $\mathsf{KeyGen}(t, n) \to (pk, \{sk_i\}_{i=1}^n)$:
        distributed key generation producing a joint public key $pk$ and
        key shares.
        \item $\mathsf{Sign}(S, m) \to \sigma$: signing protocol among
        a qualified set $S$ ($|S| \geq t$) producing signature $\sigma$
        on message $m$.
        \item $\mathsf{Verify}(pk, m, \sigma) \to \{0,1\}$: standard
        single-key verification.
    \end{enumerate}
\end{definition}

\subsection{Security Model}

\begin{definition}[Unforgeability]
    A threshold signature scheme is \emph{existentially unforgeable under
    chosen message attack} (EUF-CMA) if no PPT adversary controlling
    $t-1$ parties can produce a valid signature on a new message, even
    after observing polynomially many signing sessions.
\end{definition}

% ============================================================================
\section{FROST for Pars Guardian Networks}
\label{sec:frost}

\subsection{Protocol Overview}

FROST~\cite{komlo2020frost} is a two-round Schnorr threshold signing
protocol.  We instantiate it over Ed25519 for the Pars session layer.

\paragraph{Round 1: Commitment.}
Each signer $P_i$ in the signing set $S$ samples nonces
$(d_i, e_i) \leftarrow \mathbb{Z}_q^2$ and broadcasts commitments
$(D_i, E_i) = (d_i \cdot G, e_i \cdot G)$.

\paragraph{Round 2: Signature Share.}
Given message $m$ and all commitments, each signer computes:
\begin{align}
    \rho_i &= H_1(i, m, \{D_j, E_j\}_{j \in S}) \\
    R &= \sum_{j \in S} (D_j + \rho_j \cdot E_j) \\
    c &= H_2(R, pk, m) \\
    z_i &= d_i + \rho_i \cdot e_i + \lambda_i \cdot sk_i \cdot c
\end{align}
where $\lambda_i$ is the Lagrange coefficient for participant $i$ in
set $S$.

\paragraph{Aggregation.}
The coordinator computes $z = \sum_{i \in S} z_i$ and outputs
$(R, z)$ as the Schnorr signature.

\begin{theorem}[FROST Security]
    FROST is EUF-CMA secure under the algebraic one-more discrete
    logarithm (AOMDL) assumption in the random oracle
    model~\cite{komlo2020frost}.
\end{theorem}

\subsection{Guardian Social Recovery}

In the Pars social recovery system, a user designates $n$ guardians
with threshold $t$.  The user's master key $sk$ is shared via Feldman
VSS among the guardians.

\begin{enumerate}
    \item \textbf{Registration.}  User runs $\mathsf{KeyGen}(t, n)$ with
    guardians, distributing shares $\{sk_i\}$ and publishing $pk$.
    \item \textbf{Normal operation.}  User signs with their local key;
    guardians are passive.
    \item \textbf{Recovery.}  User contacts $t$ guardians via the Pars
    session layer (PNP-003).  Guardians run FROST to sign a key rotation
    transaction that installs the user's new key.
    \item \textbf{Verification.}  The Pars L2 contract verifies the
    threshold signature against the registered $pk$ and executes the
    key rotation.
\end{enumerate}

\subsection{Key Resharing}

Guardian sets change over time: guardians may become unreachable,
compromised, or the user may wish to adjust the threshold.

\begin{definition}[Key Resharing]
    Given an existing $(t, n)$ sharing $\{sk_i\}_{i=1}^n$ of secret $sk$
    with public key $pk$, a resharing protocol produces a new
    $(t', n')$ sharing $\{sk'_j\}_{j=1}^{n'}$ of the \emph{same} secret
    $sk$, with the same public key $pk$.
\end{definition}

The resharing protocol runs as a FROST signing session that authorises
the transition, followed by a new Feldman VSS from the old qualified
set to the new guardian set.

\begin{theorem}[Resharing Security]
    After resharing, the old shares $\{sk_i\}$ are information-theoretically
    independent of the new shares $\{sk'_j\}$.  An adversary holding
    $t-1$ old shares and $t'-1$ new shares learns nothing about $sk$.
\end{theorem}

\subsection{Proactive Refresh}

Even without guardian changes, periodic share refresh prevents gradual
compromise.  Each guardian $P_i$ shares a random polynomial of degree
$t-1$ with constant term $0$, and all guardians add the received shares
to their current shares.  The secret $sk$ and public key $pk$ remain
unchanged, but all shares are refreshed.

% ============================================================================
\section{CGGMP21 for the Diaspora Bridge}
\label{sec:cggmp}

\subsection{Motivation}

The Pars diaspora bridge connects the Pars L2 to Ethereum for
remittances.  Ethereum requires ECDSA signatures on secp256k1.
Schnorr-based FROST cannot produce valid ECDSA signatures, necessitating
a dedicated ECDSA threshold protocol.

\subsection{Protocol Overview}

CGGMP21~\cite{canetti2021cggmp} is a threshold ECDSA protocol with
identifiable abort: if a party misbehaves, the honest parties can
identify and exclude the malicious participant.

\paragraph{Key Generation.}
Parties run a distributed key generation protocol producing Paillier
key pairs for each party, a shared ECDSA public key $Q$, and additive
key shares $x_i$ such that $\sum x_i = x$ where $Q = x \cdot G$.

\paragraph{Pre-signing.}
Before the message is known, parties generate multiplicative-to-additive
(MtA) shares of $k \cdot x$ and $k \cdot \gamma$ where $k, \gamma$ are
random nonces.  This enables a single-round online phase.

\paragraph{Signing.}
Given message hash $m$:
\begin{enumerate}
    \item Each party computes partial signature $\sigma_i$.
    \item Coordinator aggregates: $\sigma = \sum \sigma_i \bmod q$.
    \item Output $(r, s)$ where $r$ is derived from $R = k^{-1} \cdot G$.
\end{enumerate}

\begin{theorem}[CGGMP21 Security]
    CGGMP21 is EUF-CMA secure under the strong RSA assumption (for
    Paillier) and the ECDSA unforgeability assumption, with
    identifiable abort~\cite{canetti2021cggmp}.
\end{theorem}

\subsection{Bridge Signing Workflow}

\begin{enumerate}
    \item A user initiates a bridge transaction on Pars L2 (lock tokens).
    \item The bridge relayer detects the lock event and constructs the
    Ethereum release transaction.
    \item The signing committee ($t$-of-$n$ bridge operators) runs
    CGGMP21 to produce the ECDSA signature.
    \item The relayer submits the signed transaction to Ethereum.
    \item Confirmation is relayed back to Pars L2.
\end{enumerate}

\subsection{Identifiable Abort}

If a bridge operator submits invalid MtA proofs or partial signatures,
the other operators can identify the malicious party and exclude them.
The protocol then restarts with the remaining $n-1$ operators (provided
$n - 1 \geq t$).

This is critical for the diaspora bridge: a state actor who compromises
one bridge operator cannot silently block all bridge transactions.

% ============================================================================
\section{Coercion-Resistant Signing}
\label{sec:coercion}

The Pars coercion resistance model (PNP-005) extends to threshold
signing.  A coerced signer must be able to participate in signing
without revealing that they are under duress.

\subsection{Panic Signing}

\begin{definition}[Panic Signing]
    A signer $P_i$ under coercion activates their panic credential
    (PNP-005).  The panic signing protocol:
    \begin{enumerate}[label=(\roman*)]
        \item Derives a decoy key share $sk_i^*$ from the panic seed.
        \item Participates in the FROST protocol using $sk_i^*$.
        \item The resulting signature is valid under a decoy public key
        $pk^*$ associated with the decoy wallet.
    \end{enumerate}
    The coercer observes a valid signing session and a valid signature
    but cannot distinguish it from a legitimate signing session.
\end{definition}

\subsection{Alert Propagation}

When a panic signing event occurs, the system must alert non-coerced
guardians without alerting the coercer:

\begin{enumerate}
    \item The panic signing session produces a valid decoy signature.
    \item Simultaneously, a steganographic alert is embedded in the
    session layer traffic (PNP-003).
    \item Non-coerced guardians detect the alert and initiate emergency
    key resharing, rotating out the compromised guardian.
\end{enumerate}

% ============================================================================
\section{Communication Layer}
\label{sec:comm}

All threshold signing rounds are conducted over the Pars Session Protocol
(PNP-003), which provides:

\begin{itemize}
    \item End-to-end encryption between signers.
    \item Forward secrecy: compromise of a signer's session keys does
    not reveal past signing transcripts.
    \item Deniability: no signer can prove to a third party that
    another signer participated.
    \item Mesh routing (PNP-002): signing sessions can proceed even
    during internet shutdowns via local mesh.
\end{itemize}

\begin{theorem}[Signing Transcript Deniability]
    Under the interactive deniability model of PNP-005, threshold signing
    transcripts conducted over PSP are deniable: no verifier can
    distinguish a real transcript from a simulated one, given only the
    final signature.
\end{theorem}

% ============================================================================
\section{Performance}
\label{sec:performance}

\begin{table}[h]
\centering
\caption{Threshold signing benchmarks (3-of-5 configuration).}
\label{tab:bench}
\begin{tabular}{@{}lrrr@{}}
\toprule
\textbf{Protocol} & \textbf{Rounds} & \textbf{Latency (ms)} & \textbf{Comm. (KB)} \\
\midrule
FROST KeyGen   & 2 & 320 & 12.4 \\
FROST Sign     & 2 & 85  & 2.1 \\
FROST Reshare  & 2 & 410 & 18.6 \\
CGGMP21 KeyGen & 4 & 2800 & 384.0 \\
CGGMP21 PreSign & 3 & 1200 & 96.0 \\
CGGMP21 Sign   & 1 & 45  & 0.8 \\
\bottomrule
\end{tabular}
\end{table}

FROST signing at 85\,ms is suitable for interactive guardian recovery.
CGGMP21 pre-signing is expensive (1.2\,s) but can be pre-computed;
online signing takes only 45\,ms.

\begin{table}[h]
\centering
\caption{Mobile device benchmarks (iPhone 15, Android Pixel 8).}
\label{tab:mobile}
\begin{tabular}{@{}lrr@{}}
\toprule
\textbf{Operation} & \textbf{iOS (ms)} & \textbf{Android (ms)} \\
\midrule
FROST Sign   & 120 & 145 \\
FROST Reshare & 580 & 710 \\
CGGMP21 Sign & 65  & 82 \\
\bottomrule
\end{tabular}
\end{table}

% ============================================================================
\section{Security Analysis}
\label{sec:security}

\subsection{Adversary Model}

We consider a static adversary who corrupts at most $t-1$ signers
before protocol execution begins.

\begin{theorem}[Unforgeability]
    An adversary corrupting $t-1$ of $n$ signers cannot produce a valid
    signature on any message not signed by the honest signers, under the
    AOMDL assumption (FROST) or strong RSA + ECDSA assumptions (CGGMP21).
\end{theorem}

\begin{theorem}[Key Secrecy]
    An adversary corrupting $t-1$ signers learns no information about the
    secret key $sk$ beyond what is implied by the public key $pk$ and
    observed signatures.
\end{theorem}

\subsection{Composition with Pars Protocols}

\begin{proposition}[UC Composition]
    When composed with the Pars session layer (PNP-003) and coercion
    resistance (PNP-005) via Universal Composability, the threshold
    signing protocols preserve their unforgeability and secrecy
    properties.
\end{proposition}

% ============================================================================
\section{Deployment}
\label{sec:deployment}

\subsection{Guardian Selection}

For social recovery, we recommend:
\begin{itemize}
    \item $n = 5$, $t = 3$ for individual accounts.
    \item $n = 7$, $t = 4$ for organisational accounts.
    \item Guardians distributed across at least 3 jurisdictions.
    \item At least one guardian reachable via mesh network (PNP-002).
\end{itemize}

\subsection{Bridge Committee}

For the diaspora bridge:
\begin{itemize}
    \item $n = 9$, $t = 6$ initial committee.
    \item Rotating membership via governance vote.
    \item Proactive refresh every 24 hours.
    \item Geographic distribution across diaspora hubs.
\end{itemize}

\subsection{Governance Execution}

For governance proposal execution:
\begin{itemize}
    \item Committee size determined by DAO governance.
    \item Threshold set to simple majority: $t = \lfloor n/2 \rfloor + 1$.
    \item Time-locked execution with emergency veto via
    $\lceil 2n/3 \rceil$ threshold.
\end{itemize}

% ============================================================================
\section{Related Work}
\label{sec:related}

\paragraph{FROST.}  Komlo and Goldberg~\cite{komlo2020frost} introduced
FROST as a two-round Schnorr threshold scheme.  We extend FROST with
key resharing and coercion-resistant panic signing.

\paragraph{CGGMP21.}  Canetti et al.~\cite{canetti2021cggmp} presented
threshold ECDSA with identifiable abort.  We deploy CGGMP21 specifically
for the Ethereum bridge, adding proactive refresh.

\paragraph{GG20.}  Gennaro and Goldfeder~\cite{gennaro2020gg20}
introduced an earlier threshold ECDSA protocol.  CGGMP21 improves on
GG20 with identifiable abort, critical for the bridge adversary model.

\paragraph{Lux Threshold Primitives.}  The Lux network provides
threshold BLS signatures for consensus.  Pars reuses the Lux threshold
infrastructure where compatible and adds Ed25519/ECDSA threshold
protocols for application-layer use.

% ============================================================================
\section{Conclusion}
\label{sec:conclusion}

We have presented a threshold signature framework for the Pars Network
that secures guardian-based social recovery, multi-party governance
execution, and diaspora bridge signing.  FROST provides efficient
Ed25519 threshold signatures for the session layer, while CGGMP21
provides ECDSA threshold signatures for Ethereum interoperability.
Key resharing enables guardian rotation, proactive refresh prevents
gradual compromise, and coercion-resistant signing extends the Pars
threat model to threshold operations.

The framework eliminates single points of failure in all critical
signing operations while maintaining sub-second latencies on consumer
hardware, enabling practical deployment across the geographically
distributed Persian diaspora.

% ============================================================================
\begin{thebibliography}{99}

\bibitem{komlo2020frost}
C.~Komlo and I.~Goldberg,
``FROST: Flexible Round-Optimized Schnorr Threshold Signatures,''
\emph{SAC 2020}, pp.~34--65, 2020.

\bibitem{canetti2021cggmp}
R.~Canetti, R.~Gennaro, S.~Goldfeder, N.~Makriyannis, and U.~Peled,
``UC Non-Interactive, Proactive, Threshold ECDSA with Identifiable
Aborts,''
\emph{CCS 2020}, pp.~1769--1787, 2021.

\bibitem{gennaro2020gg20}
R.~Gennaro and S.~Goldfeder,
``One Round Threshold ECDSA with Identifiable Abort,''
\emph{IACR ePrint 2020/540}, 2020.

\bibitem{feldman1987vss}
P.~Feldman,
``A Practical Scheme for Non-Interactive Verifiable Secret Sharing,''
\emph{FOCS}, pp.~427--437, 1987.

\bibitem{shamir1979}
A.~Shamir,
``How to Share a Secret,''
\emph{Communications of the ACM}, vol.~22, no.~11, pp.~612--613, 1979.

\bibitem{pars_social_recovery}
Cyrus Pahlavi, Pars Team,
``Social Recovery Wallets: Guardian-Based Account Recovery for Diaspora
Communities,''
2026.

\bibitem{pars_shadow_governance}
Cyrus Pahlavi, Pars Team,
``Shadow Governance: Anonymous Proposal and Execution for Diaspora DAOs,''
2026.

\bibitem{pars_diaspora_bridge}
Cyrus Pahlavi, Pars Team,
``Diaspora Bridge: Cross-Chain Remittances for the Persian Community,''
2026.

\end{thebibliography}

\end{document}
